\section{Episodic records and abstraction loss}\label{sec:C14:episodes}
An episode should preserve the relationship between the situation, the policy's attempt, the observation, and the evaluated outcome. A useful record distinguishes a successful action from a plausible explanation of why it succeeded. If the environment returned a success signal, the agent can record that signal. Inferring a general strategy from the signal is a further abstraction. Keeping those layers separate prevents one lucky trajectory from becoming an apparently validated rule.

Let a trajectory be $\xi=(o_0,a_0,o_1,\ldots,a_{T-1},o_T)$ and let $z=s(\xi)$ be its episodic abstraction. A downstream retrieval query may need an exact error message, a resource identifier, or a failed alternative that the summary omitted. The abstraction's quality depends on the future query distribution, as in Chapter~\ref{ch:C11}. A narrative that captures the overall task can still be insufficient for a later repair. Store a pointer to the source trajectory when retention permits, and preserve exact evidence fields that support verification.

RecMem's code illustrates a nontrivial failure boundary. Its consolidation routine sorts existing raw neighbors by timestamp and then appends the incoming record, rather than inserting it into that chronological order. It constructs episodic memories and proceeds to semantic refinement. If episodic generation returns no records, it returns before deleting raw neighbors. If semantic extraction returns no facts, the inspected path can still proceed to remove consolidated raw items~\cite{repo_recmem2026}. These are implementation observations at the cited revision, not measured paper failures. A backdated incoming record provides a discriminating test of prompt order versus valid time; generation failure provides a separate test of which evidence survives each storage-layer transition.

There is a narrower ingestion gap on that early-return path: the incoming record is passed directly to consolidation, and the caller does not then insert it into the raw store. Preserving existing neighbors therefore does not prove that the newly arrived interaction remains recoverable. This is visible in the inspected control flow, not an independently reproduced loss event. A fault-injection test must check the incoming item separately from the old raw neighbors, with an upstream acknowledged event log as a distinct baseline.

The corresponding engineering question is atomicity of abstraction. Should raw evidence be removed only after every intended derivative is committed, or is the episodic artifact sufficient to justify removal? There is no universal answer. If the task requires fine-grained facts that episodic extraction can omit, deletion after partial refinement may reduce recoverability. If raw content has a strict retention limit, preserving it indefinitely can violate the application contract. The correct behavior requires an explicit transaction and retention policy.

Episode merging also risks losing exception structure. Two similar successes can differ in a precondition that later determines whether a procedure applies. Merging them into a generic account may hide that condition. An episode store should therefore preserve source membership and avoid claiming that a merged summary contains every distinction from its constituents. A proposed merge validator tests representative questions and action preconditions against both the originals and the candidate abstraction, with failures retained as counterexamples.

Outcome metadata needs comparable care. A task completed under oracle feedback, a permissive simulator, or an outdated tool version can be useful experience, but its outcome should not be treated as evidence for an unchanged production environment. The episode's environment and evaluator versions are part of its meaning. Without them, retrieval can faithfully surface experience whose apparent success rests on conditions that no longer hold.

